\hypertarget{appendix-2e7e9fd28a2c}{}
\section*{Rendering history validity}
\section{Rendering history validity}\label{rendering-history-validity}

A temporal renderer reuses information from earlier frames only if that information still estimates the same target quantity at the reprojected location. ``Previous frame exists'' is therefore not enough. History needs a validity contract.

A history sample may become invalid because of:

\begin{itemize}
\tightlist
\item
  camera cuts or projection changes;
\item
  resize or dynamic-resolution changes;
\item
  object transform discontinuities;
\item
  newly visible or disoccluded surfaces;
\item
  material or lighting changes;
\item
  estimator changes such as sample pattern, exposure domain or denoiser configuration;
\item
  resource reallocation or version replacement.
\end{itemize}

\subsection{Reprojection is a hypothesis}\label{reprojection-is-a-hypothesis}

Motion vectors or previous transforms propose a correspondence from current sample to previous sample. The renderer should test that hypothesis with available evidence such as depth, normal, object/material identity, velocity magnitude or neighborhood statistics.

A common failure is using a valid coordinate that points to the wrong surface. Bounds checking prevents memory errors, not temporal ghosting.

\subsection{Versioned history}\label{versioned-history}

Treat history like a cache with a key. Relevant state changes advance a version or invalidate specific history resources. The render graph can then express that a temporal pass consumes history version N and produces version N+1.

This is also important for progressive path tracing. Accumulation belongs to a camera/scene/material/light/estimator state. If any target-defining state changes and old samples remain, the output converges to a mixture of different images.

\subsection{Failure-oriented tests}\label{failure-oriented-tests}

Useful temporal tests include a hard camera cut, a moving foreground object revealing background, an emissive material changing abruptly, resize, FOV change, jitter-sequence restart, and an object whose ID remains but geometry changes. For each case, record which history samples should survive and why.

The existing render-graph witness proves logical resource dependencies, not temporal validity itself. Native renderer evidence should capture motion/depth/ID buffers and show the acceptance/rejection mask.

Teach and reconstruct: \href{Continuous\%20Notes/Game\%20Engine\%20-\%20From\%20First\%20Frame\%20to\%20Living\%20Worlds\#Let\%20a\%20renderer\%20evolve\%20without\%20losing\%20resource\%20ownership}{Continuous Notes/Game Engine - From First Frame to Living Worlds\#Let a renderer evolve without losing resource ownership}.

Related: \hyperref[appendix-0727340a4428]{Temporal reprojection and history validation}, \href{}{Reversed-Z depth contracts (03-Visual-Computing-I, appendix-275ec8847057)}.
